\section{结论}

本文研究了长记忆 EME 信道下“离线训练 + 在线贝叶斯 CIR 适配”的神经块均衡器。离线阶段用仿真信道与完整标签训练 Pilot 条件化的展开网络；部署阶段冻结全部神经权重，仅用合法观测递推 CIR 后验。独立确认 v1 显示：四个 SNR 的配对 BER 增益均为正，H1 与 H2 通过；增益在 5/10/15 dB 随帧增长，0 dB 不通过增长判据，呈早期收敛/后段饱和。CIR 的 NMSE 下降 67\%–75\%，支持“增益来自物理状态后验”的归因。

本文的结论有明确边界：确认只运行一次，样本为 10 session × 128 帧；0 dB 的 H3 失败不能被平均掩盖；当前实现约 830 ms/帧（tail 复用约 435 ms），相对 102.4 ms 帧预算不实时；实验基于仿真信道，尚未在实测 EME 数据上验证；部署阶段网络权重保持冻结。把物理状态后验做成可部署的实时模块，并在实测信道与漂移场景中复核，是下一步工作。
